\begin{proposition}[Four-by-four images under finite-target hypotheses]
\label{B:four-by-four-finite-target-transfer}
Let $M$ be a finite monoid with a $4\times4$ rectangular-band
minimal ideal. Suppose
\[
 \dim_{\mathbb Q}H^2(BM;\mathbb Q)=1,
 \qquad H_3(BM;\mathbb Q)=0.
\]
Assume $H^2(BM;\mathbb Z)=\mathbb Zv$, that the image of $v^2$
in $H^4(BM;\mathbb Z)$ modulo torsion is nonzero and primitive,
and that every integral degree-two class restricts to zero on
every local finite subgroup of $M$.
Then the faithful row-and-column action image of $M$ has at
least twenty-four elements. In particular, $|M|\leq23$ is
impossible. No finite-CW hypothesis on $BM$ or on its action
image is required, and $H_2(BM;\mathbb Z)$ may have torsion.
\end{proposition}
\begin{proof}
We first explain why the full-hull reduction in
Proposition~\ref{B:four-by-four-CP2-hulls} still applies.
Its maximal hull is defined by the rational balanced relation
space, so the action map is an isomorphism on rational second
homology and a surjection on integral second homology.
Rank-one forms, zero fourth rational homology, and explicit
zero-square comparisons therefore give zero source squares.
Repeated rows or columns reduce to the relative three-by-four
and three-by-three arguments. Their local-character exclusions
use only the restrictions in the hypothesis: the eventual cyclic
group of a source actor surjects onto that of its image, and
inflation preserves the nonzero integral character.

The remaining even-square hulls have $H_2=\mathbb Z$, with the
displayed degree-two class its primitive dual. An integral
surjection from a finitely generated group of free rank one to
$\mathbb Z$ is an isomorphism on free parts. Its primitive dual
thus pulls back to $v$ up to sign, even if the source second
homology has torsion. The even evaluations on a saturated
integral four-cycle basis mean that the square is divisible by
two modulo torsion; the same is then true after pullback.
This contradicts the assumed nonzero primitive free square.
These are all uses of the projective-plane or finite-CW
hypotheses in that hull reduction. The image consequently lies
in one of hulls $217,216,221$.

Suppose its order is at most twenty-three. The local-character
test in Proposition~\ref{B:four-by-four-bounded22} is valid by
the same restriction hypothesis. For images of order at most
twenty-two, all its completed certificates already give zero
pulled-back rational square: every image except the thirty-nine
explicit zero-square cases has zero fourth mod-two homology.
No third-homology prefilter is used for those smaller images.

For order twenty-three, use the complete six-optional-actor
lists in Proposition~\ref{B:four-by-four-image24}. Of their
$769$ representatives with nonzero fourth mod-two homology,
four have zero third mod-two homology and are the exact
integral zero-square cases already checked there. The other
$765$ now have an independently verified rational obstruction:
\[
\begin{array}{c|r|r}
 \text{hull}&\text{representatives}&\dim_{\mathbb Q}H_3\\\hline
 217&291&1\\
 216&188&1\\
 221&286&1.
\end{array}
\]
For clarity, the new calculation is on exactly the representatives
with both $H_3(-;\mathbb F_2)$ and $H_4(-;\mathbb F_2)$ nonzero;
it does not infer rational nonvanishing from binary nonvanishing.

The normalized two-sided bar terms in degrees zero, one and
two have ranks $9,54,324$. The producer
\path{code/d_cp23_rational_h3.py} computes their exact rational
ranks and saves, for every representative, an integral cycle
$z\in\ker d_1$ and an integral functional $\ell$ with
$\ell d_2=0$ and $\ell(z)\ne0$.
The standalone checker
\path{code/d_verify_cp23_rational_h3.py} imports no producer,
classifier, or matrix-rank library. It reconstructs the literal
tables, both band actions, and both differentials, then checks
all these cycle, cocycle, and pairing identities in integers.
It verifies $9307755$ associativity equations, $563040$ band-action
equations, $2230740$ consecutive-boundary equations, and all
$765$ nonzero pairings. The complete record is
\path{code/outputs/d_cp23_rational_h3_independently_verified.json};
the integral witnesses are in
\path{code/outputs/d_cp23_rational_h3_certificates.json.gz}.

The quotient to the action image fixes the band pointwise.
Lemma~\ref{B:band-quotient-H3-surjection}, over $\mathbb Q$,
makes its third rational homology a quotient of that of $M$.
Thus none of these $765$ images can occur. Every remaining
image has zero pulled-back rational square, contrary to the
nonzero free square of $v$. This proves the assertion.
\end{proof}
